\chapter{Preface}

\section*{Why the cow}

The joke on the previous page is told against physicists, but physicists tell it about themselves, and they tell it with a certain pride. The spherical cow is not a failure of nerve. It is the whole method.

A real cow is a hopeless object. It has a nonuniform density, a complicated shape, four stomachs, and opinions. If you want to know how fast it falls off a cliff, how long it takes to cool after it dies, or how much it weighs on the Moon, almost none of that matters. A sphere of water with the right mass gets you most of the answer. The skill --- and it is a skill, the central one this book tries to teach --- is knowing \emph{which} details to throw away, throwing them away without guilt, and then knowing how to put them back when the answer stops being good enough.

You will see this move in every chapter. Near the bottom of any potential well, everything is a spring (Chapter~\ref{ch:sho}). A gas of $10^{23}$ molecules is described by three numbers (Chapter~\ref{ch:statmech}). A star is a ball of gas in balance between gravity and pressure (Chapter~\ref{ch:astro}). The universe, on large enough scales, is the same everywhere (Chapter~\ref{ch:cosmo}). Each of these is a spherical cow, and each one works astonishingly well.

\section*{What this book is}

This book is a tasting platter of physics. It was inspired by Evan Chen's \emph{An Infinitely Large Napkin}, a light but mostly self-contained tour of higher mathematics, written out of the frustration of trying to explain category theory to friends over a forty-minute lunch. The same frustration exists in physics. Someone asks what a quantum field is, and you start explaining, and you realise they need to know what a harmonic oscillator is, and then that they need to know what a Lagrangian is, and by then lunch is over.

So each chapter here is a single plate: one area of physics, served so you can taste what it is actually like. In each chapter you should get
\begin{itemize}[itemsep=2pt]
  \item the main characters and their precise definitions,
  \item one \emph{prototype} --- a single example so central that every other idea in the chapter can be tested against it,
  \item the famous results, stated precisely, with a sense of \emph{why} they are true,
  \item the orders of magnitude, so you know what the numbers actually are,
  \item a set of problems that climb from routine to genuinely hard, and
  \item directions to the real literature, so you know exactly where to go next.
\end{itemize}

There are thirty-three plates. Some are foundations that everything else rests on (mechanics, electromagnetism, quantum mechanics, statistical mechanics). Some are entire research fields compressed to a chapter (condensed matter, cosmology, particle physics). Some are crossovers where physics meets medicine, the Earth, life, or computers. Taken together they cover, at a first-taste level, roughly everything you would find in the physics section of a university library or a preprint server.

\section*{What this book is not}

It is not a replacement for a textbook. A chapter here cannot do what a semester of Griffiths or a summer with Taylor can. The amount of depth is not comparable, and it is not meant to be. Physics is learned by doing problems, many of them, over a long time, and the honest advice is still the old advice: pick up the standard books and work through them with a pencil.

What this book tries to do is make that possible. Its goal is that after any chapter you can open the serious book on that subject and \emph{recognise the landscape}: you know what the main objects are, you have met the prototype, you know which results are coming and roughly why they hold. The rest is filling in. Every chapter ends with a short guide to the literature for exactly this reason.

\section*{Three mottos}

There are three ideas that run through everything here. They are worth stating once, plainly.

\begin{keyidea}[MOTTO 1]
\textbf{Physical intuition supersedes derivations.} A derivation tells you that something is true. An explanation tells you why it had to be. Where a derivation teaches nothing, it is skipped; where it contains an idea worth seeing, it stays.
\end{keyidea}

\begin{keyidea}[MOTTO 2]
\textbf{Estimate before you calculate.} Before solving anything, guess the answer to within a factor of ten using units, limiting cases, and orders of magnitude. If the exact answer disagrees wildly with the guess, one of them is wrong, and finding out which one is how you learn.
\end{keyidea}

\begin{keyidea}[MOTTO 3]
\textbf{Learn the prototype, then test everything against it.} For each new definition or theorem, ask what it says about the prototype. If the prototype is a good one, the definition will seem natural and the theorem will seem obvious.
\end{keyidea}

\section*{Who this is for}

The book was written with a curious, mathematically capable high-school student or early undergraduate in mind: someone comfortable with calculus who wants to know what the rest of physics is actually about. If you are further along, you may still find it useful as a way to see subjects you never studied, or to remember why the subjects you did study fit together. If you are earlier along, Chapter~\ref{ch:math} tries to supply the mathematics you will need, though it moves quickly.

\section*{Acknowledgements}

This book owes its structure and spirit to Evan Chen's \emph{Napkin}, whose shape it borrows without apology, and its sense of what a solid foundation looks like to Kevin Zhou's writing on learning physics and his short list of canonical texts. Its voice owes a great deal to the \emph{Feynman Lectures on Physics}, to Purcell's \emph{Electricity and Magnetism}, and to David Morin's problem books, which are proof that problems can be a joy. Every textbook recommended at the end of a chapter is recommended because it taught somebody, somewhere, how this subject really works.
